\sectionkicker{Source-backed technical notes}
\subsection*{Agents \& Memory — 長時間動作を支える状態とスケジューリング: Technical Notes}
本欄は記事本文で圧縮した一次資料上の情報を、比較・再検証しやすい形で整理したものである。時系列、確認済みの主張、留意点、一次資料URLを掲載する。
\medskip
\subsection*{Theme at a glance}
\addcontentsline{toc}{subsection}{Theme at a glance}
\begin{center}
\footnotesize
\begin{tabularx}{\linewidth}{@{}p{0.20\linewidth}p{0.10\linewidth}p{0.14\linewidth}X@{}}
\toprule
資料 & 位置づけ & 種別 & 時系列 \\
\midrule
SwarmX: Agentic Scheduling for Low-Latency Agentic Systems & 主要資料 & 論文 & 2026-06-19 (論文\_RELEASE) \\
ChatGPT Dreaming memory architecture & 主要資料 & PRODUCT & 2026-06-04 (PRODUCT\_ROLLOUT) \\
SmoothAgent: Efficient Long-Horizon LLM-Based Agent Serving with Lookahead Context Engineering & 補足資料 & 論文 & 2026-06-30 (論文\_RELEASE) \\
\bottomrule
\end{tabularx}
\end{center}
\normalsize

\begin{technicalnote}{SwarmX: Agentic Scheduling for Low-Latency Agentic Systems}{主要資料}
\begin{tabularx}{\linewidth}{@{}>{\bfseries}p{0.22\linewidth}X@{}}
組織 & Paper authors \\
種別 & 論文 \\
時系列 & 2026-06-19 (論文\_RELEASE) \\
\end{tabularx}
\smallskip
{\bfseries 一次資料から整理したtechnical points}
\begin{itemize}[leftmargin=1.5em,itemsep=0.35em]
\item \textbf{Author claim}: The authors propose SwarmX, a scheduler designed around multi-call and tool-using agent workloads rather than independent inference requests.
\end{itemize}
{\bfseries 読む際の境界}
\begin{itemize}[leftmargin=1.5em,itemsep=0.35em]
\item \textbf{分析上の留意点}: The evidence is the authors' arXiv paper/model card and reported evaluations; results are not independently reproduced in this Evidence review.
\end{itemize}
{\bfseries 一次資料}
\begingroup\sloppy
\begin{itemize}[leftmargin=1.5em,itemsep=0.25em]
\item {\scriptsize\url{http://arxiv.org/abs/2606.21401v2}}
\end{itemize}
\endgroup
\end{technicalnote}
\medskip

\begin{technicalnote}{ChatGPT Dreaming memory architecture}{主要資料}
\begin{tabularx}{\linewidth}{@{}>{\bfseries}p{0.22\linewidth}X@{}}
組織 & OpenAI \\
種別 & PRODUCT \\
時系列 & 2026-06-04 (PRODUCT\_ROLLOUT) \\
\end{tabularx}
\smallskip
{\bfseries 一次資料から整理したtechnical points}
\begin{itemize}[leftmargin=1.5em,itemsep=0.35em]
\item \textbf{Vendor claim}: OpenAI launched a more capable and compute-efficient ChatGPT memory architecture built on dreaming, with a reviewable memory summary and staged rollout.
\end{itemize}
{\bfseries 読む際の境界}
\begin{itemize}[leftmargin=1.5em,itemsep=0.35em]
\item \textbf{分析上の留意点}: The source is first-party OpenAI material. Capability, benchmark, efficiency, or performance statements remain vendor claims unless explicitly described as externally reproduced.
\end{itemize}
{\bfseries 一次資料}
\begingroup\sloppy
\begin{itemize}[leftmargin=1.5em,itemsep=0.25em]
\item {\scriptsize\url{https://openai.com/index/chatgpt-memory-dreaming}}
\end{itemize}
\endgroup
\end{technicalnote}
\medskip

\begin{technicalnote}{SmoothAgent: Efficient Long-Horizon LLM-Based Agent Serving with Lookahead Context Engineering}{補足資料}
\begin{tabularx}{\linewidth}{@{}>{\bfseries}p{0.22\linewidth}X@{}}
組織 & Paper authors \\
種別 & 論文 \\
時系列 & 2026-06-30 (論文\_RELEASE) \\
\end{tabularx}
\smallskip
{\bfseries 一次資料から整理したtechnical points}
\begin{itemize}[leftmargin=1.5em,itemsep=0.35em]
\item \textbf{Author claim}: The authors propose SmoothAgent to execute context transformations ahead of time so long-horizon agent serving can preserve low TTFT.
\end{itemize}
{\bfseries 読む際の境界}
\begin{itemize}[leftmargin=1.5em,itemsep=0.35em]
\item \textbf{分析上の留意点}: The evidence is the authors' arXiv paper/model card and reported evaluations; results are not independently reproduced in this Evidence review.
\end{itemize}
{\bfseries 一次資料}
\begingroup\sloppy
\begin{itemize}[leftmargin=1.5em,itemsep=0.25em]
\item {\scriptsize\url{http://arxiv.org/abs/2607.00151v1}}
\end{itemize}
\endgroup
\end{technicalnote}
\medskip
